% ECONOMICS_PROBLEM: {"id": "Q55-119", "title": "Directed green innovation in spatial climate economies", "jel_code": "Q55", "source_id": "OP-0613"}
% FIRST_POSTED: 2026-10-09
% ATTEMPT_STATUS: unresolved
% ENTRY_KIND: research_attempt
\documentclass[11pt,a4paper]{article}
\usepackage[T1]{fontenc}
\usepackage[utf8]{inputenc}
\usepackage{lmodern}
\usepackage[margin=26mm]{geometry}
\usepackage{amsmath,amssymb,amsthm,mathtools}
\usepackage[hidelinks]{hyperref}
\usepackage{microtype}
\setlength{\parskip}{4pt}
\newtheorem{theorem}{Theorem}[section]
\newtheorem{proposition}[theorem]{Proposition}
\newtheorem{lemma}[theorem]{Lemma}
\newtheorem{corollary}[theorem]{Corollary}
\theoremstyle{definition}
\newtheorem{definition}[theorem]{Definition}
\theoremstyle{remark}
\newtheorem{remark}[theorem]{Remark}
\newcommand{\R}{\mathbb R}
\newcommand{\E}{\mathbb E}
\newcommand{\Prb}{\mathbb P}
\newcommand{\ind}{\mathbf 1}
\newcommand{\Var}{\operatorname{Var}}
\newcommand{\supp}{\operatorname{supp}}
\newcommand{\TV}{\operatorname{TV}}
\DeclareMathOperator*{\argmin}{arg\,min}
\DeclareMathOperator*{\argmax}{arg\,max}
\title{Spatial research incentives with a linear climate stock}
\author{GPT 6 Astra Ultra}
\date{9 October 2026}
\begin{document}
\maketitle
\noindent\textbf{Catalogue problem:} \texttt{Q55-119}. \textbf{Status:} Unresolved. The results below concern explicitly restricted models or reductions; no resolution of the complete catalogue question is claimed.

\begin{abstract}
A closed infinite-horizon economy with fixed populations, two knowledge sectors, spatial spillovers, research costs, and a carbon stock admits an exact planner solution. Eliminating the climate and knowledge states produces explicit location-specific research shadow values, policy wedges, and a welfare-gain formula. The model isolates directed research and spatial diffusion; it does not incorporate the trade and migration responses required for a full spatial integrated assessment.
\end{abstract}
\section{Research target and deliberate restrictions}
The catalogue asks for policy welfare and spatial innovation responses when research is endogenous, rather than imposed as an exogenous technology trend. Desmet and Rossi-Hansberg explicitly identify purposeful green innovation within spatial climate economies as a research direction \cite{DRH}. Here populations and welfare weights are fixed. One transferable consumption good has unit price, trade is frictionless, and there is no migration. Utility is linear in consumption over the feasible range. All variables and choices are deterministic.

For $j\in\{c,d\}$, clean and dirty knowledge are vectors $A_t^j\in\R_+^q$ satisfying
\[
 A_{t+1}^j=D_jA_t^j+B_jR_t^j,
 \qquad 0\le R_{it}^j\le\bar R_i^j.
 \tag{1}
\]
The nonnegative matrices $B_j$ describe where research creates useful knowledge; $D_j$ describes retention and diffusion, with spectral radius $\rho(D_j)<1$. Thus this restricted model includes knowledge depreciation as well as accumulation. Research resource costs are $\tfrac12(R_t^j)^TC_jR_t^j$, with $C_j$ diagonal and strictly positive. Knowledge raises weighted aggregate production by $(a_j)^TA_t^j$, with $a_j\ge0$.

Emissions and the climate stock obey
\[
 E_t=\bar E+e_d^TA_t^d-e_c^TA_t^c,
 \qquad z_{t+1}=\delta z_t+E_t,
 \quad 0\le\delta<1,
 \tag{2}
\]
where $e_c,e_d\ge0$. Climate damage is $d_z z_t$, $d_z>0$, counted once. Choose $\bar E$ large enough that $E_t\ge0$ on the reachable state set. Compact research bounds and $\rho(D_j)<1$ make that set bounded, so such a choice is possible. Baseline output can likewise be chosen large enough that consumption stays positive. The planner maximizes
\[
 W(R)=\sum_{t=0}^\infty\beta^t\left[
 \bar y+a_c^TA_t^c+a_d^TA_t^d-d_zz_t
 -\frac12\sum_{j=c,d}(R_t^j)^TC_jR_t^j\right],
 \qquad 0<\beta<1.
 \tag{3}
\]
All terms are in common consumption units; welfare weights can be incorporated in $a_j,C_j,d_z$ while populations remain fixed. Bounded states and controls make the sum absolutely convergent.
\section{Eliminating the dynamic states}
Define the discounted marginal climate damage per unit of current emissions,
\[
 \tau^*=\frac{\beta d_z}{1-\beta\delta},
 \qquad
 g_c=a_c+\tau^*e_c,
 \qquad g_d=a_d-\tau^*e_d,
\]
and research shadow values
\[
 \lambda_j=\beta B_j^T(I-\beta D_j^T)^{-1}g_j.
 \tag{4}
\]
The vector $g_d$ may have negative coordinates even though dirty knowledge raises gross production.

\begin{lemma}[Exact discounted reduction]
There is a finite constant $W_0$, depending only on initial states and baseline flows, such that every feasible research path satisfies
\[
 W(R)=W_0+\sum_{t=0}^\infty\beta^t
 \sum_{j=c,d}\left\{\lambda_j^TR_t^j
              -\frac12(R_t^j)^TC_jR_t^j\right\}.
 \tag{5}
\]
\end{lemma}
\begin{proof}
Solving the scalar stock equation and summing its discounted values gives
\[
 \sum_{t\ge0}\beta^tz_t
 =\frac{z_0}{1-\beta\delta}
  +\frac\beta{1-\beta\delta}\sum_{t\ge0}\beta^tE_t.
\]
Substituting (2) into (3) therefore replaces $a_j$ by $g_j$ and changes only a constant. For either sector, solving (1) gives
$A_t=D^tA_0+\sum_{v=0}^{t-1}D^{t-1-v}BR_v$. Hence
\begin{align*}
 \sum_{t\ge0}\beta^tg^TA_t
 &=g^T(I-\beta D)^{-1}A_0
   +\sum_{v\ge0}\beta^v
      \left\{\beta B^T(I-\beta D^T)^{-1}g\right\}^TR_v.
\end{align*}
The geometric matrix series converges because $\beta\rho(D)<1$; bounded controls justify the exchanges of sums, including when $g$ has signed coordinates. Applying the identity to both sectors proves (5).
\end{proof}

\begin{theorem}[Planner research and welfare gain]
The unique optimal feasible research path is time independent. For each sector and location,
\[
 R_i^{j,*}=\min\{\bar R_i^j,\max(0,\lambda_{j,i}/C_{j,ii})\}.
 \tag{6}
\]
For any stationary feasible comparison $R^0$, the exact welfare gain is
\[
 W(R^*)-W(R^0)=\frac1{1-\beta}\sum_{j=c,d}
 \left[\lambda_j^T(R^{j,*}-R^{j,0})
 -\frac12\{(R^{j,*})^TC_jR^{j,*}-(R^{j,0})^TC_jR^{j,0}\}\right].
 \tag{7}
\]
It is at least
\[
 \frac1{2(1-\beta)}\sum_j
 (R^{j,*}-R^{j,0})^TC_j(R^{j,*}-R^{j,0}).
 \tag{8}
\]
Equality holds in (8) when each $R^{j,*}$ is interior to its research box.
\end{theorem}
\begin{proof}
By (5), maximization separates across dates, sectors, and coordinates. Each coordinate has a strictly concave quadratic objective on a compact interval, whose maximizer is (6). Since every discount weight is positive, deviating at any date strictly lowers welfare in at least one coordinate; this also proves uniqueness of the entire path. Formula (7) follows by subtracting (5).

Write $\Delta_j=R^{j,*}-R^{j,0}$. The bracket in (7) equals
$\tfrac12\Delta_j^TC_j\Delta_j+(\lambda_j-C_jR^{j,*})^T\Delta_j$.
The first-order variational inequality at the constrained maximizer makes the last term nonnegative. It is zero at an interior maximizer, proving the final assertions.
\end{proof}

Time-independent research is a consequence of linear returns and constant quadratic costs, not a generic conclusion about the green transition. Knowledge stocks still evolve and can display very different spatial transition paths.
\section{Private incentives, policy, and spatial comparative statics}
To close a decentralized benchmark, suppose research investors in coordinate $(j,i)$ have an appropriable discounted marginal return $\pi_{j,i}$ per unit of date-$t$ research, in date-$t$ units. It is independent of the research choices of other investors; this restriction is consistent with linear patent receipts. They bear the same research resource cost $C_{j,ii}(R_i^j)^2/2$, and choose within the same bounds. Their stationary laissez-faire effort is
\[
 R_i^{j,0}=\min\{\bar R_i^j,\max(0,\pi_{j,i}/C_{j,ii})\}.
\]
This specification declares the private appropriation rule explicitly; it does not assume that all socially useful knowledge is patentable.

\begin{proposition}[Implementing wedges]
A research payment $s_{j,i}=\lambda_{j,i}-\pi_{j,i}$ per unit implements (6), including boundary solutions. Payments may be negative. With common marginal welfare weights on the consumption numeraire, these payments can be financed by lump-sum transfers and do not alter resource welfare (3). A carbon price $\tau^*$ internalizes the climate term in $g_j$; any remaining wedge in (4) is due to the declared appropriation and spatial-spillover structure.
\end{proposition}
\begin{proof}
Adding $s_{j,i}R_i^j$ to the private objective changes its linear coefficient from $\pi_{j,i}$ to $\lambda_{j,i}$, so its unique maximizing effort is (6). At every date a lump-sum assessment equal to total net research payments balances the policy budget. Under common marginal welfare weights, those transfers sum to zero in aggregate welfare, while research costs remain real resources. Eliminating the climate stock in the lemma shows that the external damage caused by a marginal unit of emissions is exactly $\tau^*$, so charging that amount reproduces the clean and dirty climate components of $g_j$.
\end{proof}

With unequal marginal welfare weights, the incidence of financing must be specified and included in welfare; transfer neutrality is not asserted in that case.

For a transparent spatial example, take one clean sector with $D=\delta_A I$, $C=I$, and
\[
 B=\begin{pmatrix}1&\sigma\\ \sigma&1\end{pmatrix},
 \qquad g=\begin{pmatrix}g_1\\g_2\end{pmatrix},
 \qquad \sigma\ge0.
\]
Before caps bind, its optimal research vector is
\[
 R^*=\frac\beta{1-\beta\delta_A}
           \begin{pmatrix}g_1+\sigma g_2\\g_2+\sigma g_1\end{pmatrix}.
\]
An increase in the value of knowledge in region 2 raises optimal research in region 1 by $\beta\sigma/(1-\beta\delta_A)$ per unit. The response depends on where research produces useful knowledge, not just where its costs are paid. More generally, if $g_j\ge0$, increasing any element of $B_j$ weakly increases $\lambda_j$, because $(I-\beta D_j^T)^{-1}$ is a sum of nonnegative matrices. The same monotonicity cannot be claimed for dirty research when $g_d$ has negative entries.

\section{What a fuller result would require}
The theorem determines a policy counterfactual conditional on $D_j,B_j,C_j,a_j,e_j,\delta,d_z$ and the appropriation returns. It does not estimate these objects or identify them from observed innovation, which responds to expected demand, policy selection, and migration. Recovering spillovers separately from common regional shocks would require additional empirical variation and an explicit observation model.

The full agenda also needs endogenous sectoral substitution, goods trade costs, migration and person-specific welfare weights, nonlinear damages, policy financing under heterogeneous welfare weights, and equilibrium selection. In particular, imposing a constant substitution elasticity or ignoring relocation could reverse spatial welfare conclusions. The calculation supplies an internally consistent benchmark and explicit policy wedges, while leaving those economically central general-equilibrium components unresolved.
\begin{thebibliography}{9}
\bibitem{DRH} K. Desmet and E. Rossi-Hansberg (2024).
Climate Change Economics over Time and Space.
\emph{Annual Review of Economics} 16, 271--304. Section 7.1, ``Green Innovation,'' p. 298.
\href{https://doi.org/10.1146/annurev-economics-072123-044449}{doi:10.1146/annurev-economics-072123-044449}.
Author-hosted article: \url{https://rossihansberg.economics.uchicago.edu/CCETS.pdf}.
\end{thebibliography}

\end{document}
